\section{Выводы}
{\itshape Выполнив лабораторную работу по курсу <<Дискретный анализ>>, я изучил и реализовал 
битовую версию структуры данных PATRICIA, а также закрепил навыки работы с низкоуровневыми 
представлениями данных.}

В ходе работы я пришел к следующим выводам:
\begin{itemize}
    \item \textbf{Эффективность bitwise PATRICIA:} Использование побитового сравнения позволяет 
    минимизировать количество операций сравнения и ускорить выполнение поиска, вставки и удаления.
    
    \item \textbf{Отличие от символьных деревьев:} В отличие от классических префиксных деревьев, 
    где сравнение происходит посимвольно, в PATRICIA сравнение выполняется на уровне битов, 
    что делает структуру более компактной и быстрой.
    
    \item \textbf{Работа с указателями:} Реализация потребовала аккуратного управления указателями 
    и понимания структуры дерева, особенно при реализации операции удаления.
    
    \item \textbf{Бинарное представление данных:} Работа с битами и индексами битов позволила лучше 
    понять внутреннее устройство строк и способы их эффективной обработки.
    
    \item \textbf{Практическая производительность:} Эксперименты показали, что PATRICIA превосходит 
    \texttt{std::map} по времени выполнения операций на больших объёмах данных.
\end{itemize}

Данный опыт полезен при разработке высокопроизводительных систем, работающих с текстовыми ключами, 
а также при проектировании специализированных структур данных.